\begin{theorem}
\label{thm:main}
Let $p\colon\moon\to[-1,1]$ be a given polynomial.\ Then
for every $\delta>0,$ there is a polynomial $p_\text{\rm
robust}\colon\R^n\to\R$ of
degree $O(\deg p+\log 1/\delta)$ such that 
\begin{align*}
   \left|p(x) - p_\text{\rm robust}(x+\epsilon)\right| <\delta
\end{align*}
for all $x\in\moon$ and $\epsilon\in[-1/3,1/3]^n$.
Furthermore, $p_\text{\rm robust}$ has an explicit, closed-form
description.
\end{theorem}

\section{Notation and preliminaries}
\label{sec:prelim}

Throughout this manuscript, we represent the Boolean
values ``true'' and ``false'' by $-1$ and $+1,$
respectively. In particular, Boolean functions are
mappings $f\colon X\to\moo$ for some finite set $X$
such as $X=\moon$. The natural numbers are denoted
$\mathbb N=\{0,1,2,3,\dots\}$. The symbol $\log x$
denotes the logarithm of $x$ to base~$2$.
For a string $x\in\R^n$ and a set $S\subseteq\oneton,$
we adopt the shorthand
$x|_S=(x_{i_1},x_{i_2},\dots,x_{i_{|S|}})\in\R^{|S|},$
where $i_1<i_2<\cdots<i_{|S|}$ are the elements of $S$. 
The family of all subsets of a given set $X$ is denoted
$\mathcal{P}(X)$. The symbol $S_n$ stands for the group
of permutations $\sigma\colon\oneton\to\oneton$. A
function
$\phi\colon\R^n\to\R$ is called \emph{symmetric}
if $\phi$ is invariant under permutations of the
variables, \ie, $\phi(x)\equiv
\phi(x_{\sigma(1)},x_{\sigma(2)},\dots,x_{\sigma(n)})$
for all $\sigma\in S_n$. 
We adopt the standard definition
of the sign function:
\begin{align*}
  \sign t = 
  \ccases{
  -1 &\text{if } t<0, \\
  0  &\text{if } t=0, \\
  1 &\text{if } t>0.
  }
\end{align*}
For a set $X,$ we let $\R^X$ denote the real vector
space of functions $X\to\R$. For $\phi\in\R^X,$ 
we write
\[
\|\phi\|_\infty =\sup_{x\in X} |\phi(x)|\,, 
\qquad\qquad \|\phi\|_1=\sum_{x\in X} |\phi(x)|\,,
\]
where the symbol $\|\phi\|_1$ is reserved for finite
$X$. By the \emph{degree} of a multivariate polynomial $p$
on $\R^n,$ denoted $\deg p,$ we shall always mean the
total degree of $p$, \ie, the greatest total degree of
any monomial of $p$. The symbol $P_d$ stands for the
family of all univariate real polynomials of degree up
to $d$. 
For a natural number $i$ and a real $\alpha,$ the
generalized binomial coefficient ${\alpha\choose i}$ is
given by \[{\alpha\choose
i}=\frac{\alpha(\alpha-1)\cdots(\alpha-i+1)}{i!}\,.\]
The following identity is well-known and straightforward to
show~\cite[p.~186]{concrete-math}:
\begin{align}
{-1/2\choose i}=\left(-\frac14\right)^i{2i\choose i}\,.
\label{eqn:gen-binom}
\end{align}

 
\subsubsection*{Fourier transform}

Consider the vector space of functions
$\moon\to\R$. For $S\subseteq\oneton,$ define
$\chi_S\colon\moon\to\moo$ by $\chi_S(x) =\prod_{i\in
S} x_i$. Then the functions $\chi_S,$
$S\subseteq\oneton,$ form an orthogonal basis for the
vector space in question.  In particular, every
function $\phi\colon\moon\to\R$ has a unique
representation as a linear combination of the
characters $\chi_S$: 
\[\phi=\sum_{S\subseteq\oneton} \hat \phi(S)\,\chi_S\,,\]
where $\hat\phi(S)=2^{-n}\sum_{x\in\moon}
\phi(x)\chi_S(x)$ is the \emph{Fourier coefficient} of
$\phi$ that corresponds to the character $\chi_S$. 
Note that 
\begin{align*}
\deg \phi= \max\{|S|:\hat \phi(S)\ne 0\}\,.
\end{align*}
Formally, the Fourier transform is the linear
transformation $\phi\mapsto\hat\phi,$ where $\hat\phi$ 
is viewed as a function 
$\mathcal{P}(\oneton)\to\R$. In particular, we have
the shorthand
\begin{align*}
\|\hat\phi\|_1=\sum_{S\subseteq\oneton}|\hat\phi(S)|\,.
\end{align*}

\subsubsection*{Multilinear extensions and convexity}

As the previous paragraph shows, associated to every mapping 
$\phi\colon\moon\to\R$ is a unique
multilinear polynomial $\tilde
\phi\colon\R^n\to\R$ 
such that $\phi\equiv \tilde \phi$ on
$\moon$. In discussing the Fourier transform, we
identified $\phi$ with its
multilinear extension $\tilde \phi$ to $\R^n,$ and 
will continue to do so throughout this paper. Among
other things, this
convention allows one to evaluate $\phi$ everywhere
in $[-1,1]^n$ as opposed to just $\moon$. 
It is a simple but important fact that for every
$\phi\colon\moon\to\R,$ 
\begin{align*}
\max_{x\in[-1,1]^n}|\phi(x)|
=\max_{x\in\moon}|\phi(x)| = \|\phi\|_\infty\,.
\end{align*}
To see this, fix $\xi\in[-1,1]^n$ arbitrarily and
consider the probability distribution on strings $x\in\moon$
whereby $x_1,\dots,x_n$ are distributed independently and
$\Exp[x_i]=\xi_i$  for all $i$. Then
$\phi(\xi)=\Exp[\phi(x)]$ by multilinearity, so that
$|\phi(\xi)|\leq\max_{x\in\moon}|\phi(x)|$.

\section{A robust polynomial for parity}
\label{sec:robust-parity}

The objective of this section is to construct
a low-degree robust polynomial
for the parity function. In 
other words, we will construct a polynomial
$p\colon\R^n\to\R$ of degree $O(n)$ such that
$p(x_1,x_2,\dots,x_n) \approx \prod \sign x_i$ whenever the
input variables are close to Boolean:
$x_1,x_2,\dots,x_n\in[-4/3,-2/3]\cup[2/3,4/3]$. 
Recall that our eventual goal is a robust polynomial
for every bounded real function. To this end, the parity
approximant $p$ that we are to construct needs to 
possess a key additional
property: the error \mbox{$p(x)-\prod\sign x_i,$} apart from
being small, needs to be expressible as a
multivariate series in which the coefficients decay
rapidly with monomial order. 

To obtain this coefficient behavior, we use a carefully
chosen approximant for the univariate function
$\sign t$. The simplest candidate is 
the following ingenious construction due to 
Buhrman et al.~\cite{robust}:
\begin{align*}
B_d(t) = 2^{-d} \sum_{i=\lceil d/2\rceil}^d {d\choose i}
t^i(1-t)^{d-i}\,.
\end{align*}
In words, $B_d(t)$ is the probability of observing
more heads than tails in a sequence of $d$ independent
coin flips, each coming up heads with probability $t$.
By the Chernoff bound, $B_d$ sends
$[0,1/4]\to[0,2^{-\Omega(d)}]$ and similarly
$[3/4,1]\to[1-2^{-\Omega(d)},1]$. As Buhrman et
al.~\cite{robust} point out, this immediately gives a
degree-$d$ approximant for the sign function with
exponentially small error on
$[-4/3,-2/3]\cup[2/3,4/3]$. Unfortunately, the
coefficients of $B_d$ do not
exhibit the rapid decay that we require.
Instead, in what follows
we use a purely analytic construction based on the
Maclaurin series for $1/\sqrt{1+t}$. 

\begin{lemma}
\label{lem:parity-robust}
For $x_1,x_2,\dots,x_n\in(-\sqrt 2,0)\cup(0,\sqrt 2),$
\begin{align}
\label{eqn:parity-maclaurin}
\sign (x_1x_2\cdots x_n) =
x_1x_2\cdots x_n \sum_{i_1,i_2,\dots,i_n\in\mathbb N} 
\;\; \prod_{j=1}^n \PARENS{-\frac14}^{i_j}{2i_j \choose i_j}
(x_j^2-1)^{i_j}\,.
\end{align}
\end{lemma}


\begin{proof}
Recall the binomial series 
$(1+t)^{\alpha} = \sum_{i=0}^\infty {\alpha \choose i}
t^i,$
valid for all $-1<t<1$
and all real $\alpha$. In particular, setting
$\alpha=-1/2$ gives
\begin{align}
\frac{1}{\sqrt{1+t}} 
   &= \sum_{i=0}^\infty {-1/2 \choose i} t^i \nonumber \\
   &= \sum_{i=0}^\infty \PARENS{-\frac14}^i {2i\choose i} t^i\,, 
   		&& -1<t<1\,,
\label{eqn:inverse-sqrt}
\end{align}
where the second step uses \eqref{eqn:gen-binom}. One easily verifies that this absolutely 
convergent series is the Maclaurin
expansion for $1/\sqrt{1+t}$. For all real $t$ with
$0<|t|<\sqrt 2,$ we have
\begin{align}
 \sign t &= \frac t{\sqrt{1 + (t^2-1)}}
  = t\sum_{i=0}^\infty \PARENS{-\frac14}^i {2i \choose i} (t^2-1)^i\,,
  \label{eqn:sign-approx-maclaurin}
\end{align}
where the second step holds by
\eqref{eqn:inverse-sqrt}. For 
$x_1,x_2,\dots,x_n\in(-\sqrt 2,0)\cup(0,\sqrt 2),$ 
it follows that
\begin{align*}
\sign(x_1x_2\dots x_n)
  &= x_1x_2\cdots x_n \prod_{j=1}^n \left\{\,\sum_{i=0}^\infty 
       \PARENS{-\frac14}^i{2i \choose i} (x_j^2-1)^i
	   \,\right\}  \\
  &= x_1x_2\cdots x_n \sum_{i_1,i_2,\dots,i_n\in\mathbb N}
  \;\; \prod_{j=1}^n
  \PARENS{-\frac14}^{i_j}{2i_j\choose i_j}(x_j^2-1)^{i_j}. 
  \qedhere
\end{align*}
\end{proof}


We have reached the main result of this section.

\begin{theorem}[Robust polynomial for the parity function]
\label{thm:parity-robust}
Fix $\epsilon\in[0,1)$ and let
\begin{align*}
X=[-\sqrt{1+\epsilon},-\sqrt{1-\epsilon}]\cup
         [\sqrt{1-\epsilon},\sqrt{1+\epsilon}]\,. 
\end{align*}
Then for every integer $N\geq1,$ there is an
$($explicitly given$)$
polynomial $p\colon\R^n\to\R$ of degree at most
$2N+n$ such that 
\begin{align}
   \max_{X^n}|\sign(x_1x_2\cdots x_n)-p(x)| \leq
   \epsilon^N(1+\epsilon)^{n/2}{N+n\choose N}n\,.
   \label{eqn:parity-error-bound}
\end{align}
\end{theorem}

Setting $\epsilon=7/9$ and $N=22n$ in this result, we infer that
the function $\sign (x_1x_2\cdots x_n)$ with inputs
$x_1,x_2,\dots,x_n\in[-4/3,-2/3]\cup[2/3,4/3]$ can be
approximated to within $2^{-n}$ everywhere 
by a polynomial of degree $O(n)$. This is the
desired robust polynomial for parity.

\begin{proof}[Proof of \expref{Theorem}{thm:parity-robust}]
If $\epsilon>N/{(N+n)},$ the right-hand side of 
\eqref{eqn:parity-error-bound} exceeds 
\[
\left(\frac N{N+n}\right)^N{N+n\choose N}\geq1
\]
and therefore the theorem holds trivially with $p=0$. 
In what follows, we treat the complementary case:
\begin{align}
\label{eqn:eps-small}
0\leq\epsilon\leq\frac N{N+n}\,.
\end{align}
For a natural number $d,$ let
$\mathcal{I}_d$ stand for the family of $n$-tuples
$(i_1,\dots,i_n)$ of nonnegative integers such that
$i_1+\cdots+i_n=d$. Clearly,
\begin{align*}
\abs{\mathcal{I}_d} = {d+n-1\choose d}\,.
\end{align*}
One can restate \eqref{eqn:parity-maclaurin} in the
form
\begin{align}
\sign(x_1x_2\cdots x_n) 
  &= x_1x_2\cdots x_n \sum_{d=0}^\infty
  \xi_d(x_1,x_2,\dots,x_n)\,,
\label{eqn:xi-series}
\end{align}
where
\begin{align*}
 \xi_d(x_1,x_2,\dots,x_n) = \sum_{(i_1,\dots,i_n)\in
 \mathcal I_d}
  \;\prod_{j=1}^n 
       \PARENS{-\frac14}^{i_j}{2i_j \choose i_j} (x_j^2-1)^{i_j}\,.
\end{align*}
On $X^n,$ 
\begin{align*}
\|\xi_d\|_\infty \leq \epsilon^d |\mathcal I_d|
=\epsilon^d{d+n-1\choose d}\,.
\end{align*}
As a result, dropping the terms
$\xi_{N+1},\xi_{N+2},\dots$ from the series
\eqref{eqn:xi-series} results in an approximant 
of degree $2N+n$ with pointwise error at most
\begin{align*}
(1+\epsilon)^{n/2}\sum_{d=N+1}^\infty
\epsilon^d{d+n-1\choose d}
&\leq (1+\epsilon)^{n/2}\,\epsilon^{N+1}{N+n\choose N+1}
\sum_{d=0}^\infty
\PARENS{\epsilon\cdot\frac{N+n}{N+1}}^d\\
&\leq (1+\epsilon)^{n/2}\,\epsilon^{N+1}{N+n\choose N+1}
\sum_{d=0}^\infty
\PARENS{\frac N{N+1}}^d 
&&\text{by \eqref{eqn:eps-small}}\\
&= (1+\epsilon)^{n/2}\,\epsilon^{N+1}{N+n\choose N}n\,.
&&\qedhere
\end{align*} 
\end{proof}
 

\section{Reduction to homogeneous polynomials}
\label{sec:homogeneous}

We now turn to the construction of a robust polynomial
for any real function on the Boolean cube. Real
functions given by homogeneous polynomials on $\moon$ are
particularly convenient to work with, and the proof is
greatly simplified by first reducing the problem to the
homogeneous case. 

To obtain this reduction, we need a bound on the
coefficients of a \emph{univariate} polynomial in terms
of its degree $d$ and its maximum absolute value on
the interval $[-1,1]$. This fundamental problem was
solved in the nineteenth century by
V.~A.~Markov~\cite[p.~81]{markov-v-a1982least-dev}, who
proved an upper bound of 
\begin{align}
O\left(\frac{(1+\sqrt 2)^d}{\sqrt d}\right)
\label{eqn:markov}
\end{align}
on the size of the coefficients of any degree-$d$
polynomial that is bounded on $[-1,1]$ in absolute value by~$1$. 
Markov further showed that \eqref{eqn:markov} is
tight. Rather than appeal to this deep result in
approximation theory, we give a first-principles proof
of a $4^d$ bound which suffices for our purposes.


\begin{lemma}[Coefficients of bounded polynomials]
\label{lem:bound-poly-coeffs}
Let $p(t)=\sum_{i=0}^d a_it^i$ be a given polynomial. Then
\begin{align}
\label{eqn:bound-poly-coeffs}
\sum_{i=0}^d \abs{a_i} \leq 4^d 
\max_{j=0,1,\dots,d}\ABS{p\PARENS{\frac {d-2j}d}}\,.
\end{align}
\end{lemma}

\begin{proof}
We will use a common approximation-theoretic 
technique~\cite{cheney-book,rivlin-book}
whereby one expresses $p$ as a linear
combination of more structured polynomials
and analyzes the latter objects. 
For this, define $q_0,q_1,\dots,q_d\in P_d$ by 
\begin{align}
 q_j(t) = \frac{(-1)^j d^d}{d!\,2^d}{d\choose j}
\prod_{\substack{i=0\\i\ne j}}^d \PARENS{t - \frac{d-2i}d}\,,
\qquad j=0,1,\dots,d\,.
\label{eqn:q-def}
\end{align}
One easily verifies that these polynomials behave like
delta functions, in the sense that for 
$i,j=0,1,2,\dots,d,$ 
\begin{align*}
q_j\PARENS{\frac {d-2i}d} = \begin{cases}
1 &\text{if $i=j,$}\\
0 &\text{otherwise.}
\end{cases}
\end{align*}
Lagrange interpolation gives
\begin{align}
p=\sum_{j=0}^d p\PARENS{\frac {d-2j}d}q_j\,.
\label{eqn:p-lin-comb-q}
\end{align}
By \eqref{eqn:q-def}, the absolute values of the
coefficients of $q_j$ sum to at most
\begin{align*}
\frac{d^d}{d!\,2^d}{d\choose j}\prod_{\substack{i=0\\i\ne
j}}^d\left(1+\frac {|d-2i|}d\right)&=
\frac{1}{d!\,2^d}{d\choose j}\cdot \frac1{d+|d-2j|}\;
\prod_{i=0}^d(d+|d-2i|)\\
&=
\frac{1}{d!\,2^d}{d\choose j}\cdot \frac1{d+|d-2j|}\;
\prod_{i=0}^{\lfloor d/2\rfloor} (2d-2i) 
\prod_{i=\lfloor d/2\rfloor+1}^d (2i) \\
&=
\frac{1}{d!\,2^d}{d\choose j}\cdot \frac{2^{d+1}}{d+|d-2j|}
\cdot \frac{d!}{\left(\lceil d/2\rceil-1\right)!} 
\cdot \frac{d!}{\lfloor d/2\rfloor!} \\
&=
{d\choose j}\cdot \frac{2d}{d+|d-2j|} {d-1 \choose \lfloor d/2\rfloor} \\
&\leq 2^d {d\choose j}\,. \end{align*}
In view of \eqref{eqn:p-lin-comb-q}, we conclude that
the absolute values of the coefficients of
$p$ sum to at most
\begin{align*}
\PARENS{\max_{j=0,1,\dots,d}\ABS{p\PARENS{\frac
{d-2j}d}}}
\sum_{j=0}^d 2^d{d\choose j} =4^d 
\max_{j=0,1,\dots,d}\ABS{p\PARENS{\frac {d-2j}d}}\,.& \qedhere
\end{align*}
\end{proof}

We are now prepared to give the desired reduction to
the homogeneous case. 

\begin{theorem}
\label{thm:rebuild-from-levels}
Let $\phi\colon\moon\to\R$ be a given function, $\deg \phi=d$.
Write $\phi=\phi_0+\phi_1+\cdots+\phi_d,$ where 
$\phi_i\colon\moon\to\R$ is given
by $\phi_i=\sum_{|S|=i}\hat\phi(S)\chi_S$. Then
\begin{align*}
\|\phi_i\|_\infty \leq 4^d \|\phi\|_\infty\,, \qquad
i=0,1,\dots,d\,.
\end{align*}
\end{theorem}

The above result gives an upper bound on the infinity
norm of the homogeneous parts of a polynomial $\phi$ in
terms of the infinity norm of $\phi$ itself.  Note that the bound
is entirely independent of the number of variables. 
For our purposes, \expref{Theorem}{thm:rebuild-from-levels}
has the following consequence: a robust 
polynomial for $\phi$ can be obtained
by constructing robust polynomials with 
error $2^{-\Omega(d)}\|\phi\|_\infty$ 
separately for each of the homogeneous parts. The
homogeneous problem will be studied in the next section.


\begin{proof}[Proof of \expref{Theorem}{thm:rebuild-from-levels}]
Pick a point $x\in\moon$ arbitrarily and fix it for the
remainder of the proof. Consider the univariate
polynomial $p\in P_d$ given by
\begin{align*}
p(t) = \sum_{i=0}^d \phi_i(x)t^i\,.
\end{align*}
For $-1\leq t\leq1,$ consider the probability
distribution $\mu_t$ on the Boolean cube $\moon$
whereby each bit is independent and has expected value
$t$. Then 
\begin{align*}
\|\phi\|_\infty
&\geq
\left\lvert\Exp_{z\sim\mu_t}[\phi(x_1z_1,\dots,x_nz_n)]\right\rvert
=
\left\lvert
\sum_{|S|\leq d} \hat\phi(S)
\Exp_{z\sim\mu_t}\left[\,\prod_{i\in S} x_iz_i\right]\right\rvert
=
\left\lvert
\sum_{|S|\leq d} \hat\phi(S)
t^{|S|} \prod_{i\in S} x_i\right\rvert
=
|p(t)|\,.
\end{align*}
Hence, $p$ is bounded on $[-1,1]$  in absolute value by
$\|\phi\|_\infty$. By
\expref{Lemma}{lem:bound-poly-coeffs}, it follows that the
coefficients of $p$ do not exceed
$4^d\|\phi\|_\infty$:
\begin{align*}
|\phi_i(x)|\leq 4^d\|\phi\|_\infty\,.
\end{align*}
Since the choice of $x\in\moon$ was arbitrary, the
theorem follows.
\end{proof}




\section{Error cancellation in homogeneous polynomials}

Recall that
the goal of this paper is to construct a degree-$O(d)$ robust
polynomial for every degree-$d$ real function
$\phi\colon\moon\to[-1,1]$. By
the results of \expref{Section}{sec:homogeneous}, we may
now assume that $\phi$ is homogeneous:
\begin{align*}
\phi = \sum_{|S|=d} \hat\phi(S)\chi_S\,.
\end{align*}
A na\"ive approach would be to use the construction of
\expref{Section}{sec:robust-parity} and robustly
approximate each parity $\chi_S$ to within
$2^{-\Theta(d)}$ by a degree-$O(d)$ polynomial.
Unfortunately, it is unclear whether the resulting
polynomial would be a good approximant for $\phi$.
Indeed, as explained in the introduction,
the cumulative error 
could conceivably be as large as $n^{\Omega(d)} 
2^{-\Theta(d)}\gg1$.
The purpose of this section is to prove that, for a
careful choice of approximants for the $\chi_S,$ the
errors do not compound but instead partially cancel,
resulting in a cumulative error of
$2^{-\Theta(d)}$. The proof is rather technical.
To simplify the exposition, we first illustrate our
technique in the simpler setting of $\moon$ and
then adapt it to our setting of interest,~$\R^n$. 

\subsection*{Error cancellation on the Boolean hypercube}
Let $\phi\colon\moon\to\R$ be a degree-$d$ homogeneous
polynomial. Our
goal is to show that perturbing the Fourier characters of
$\phi$ in a suitable, coordinated manner results in
partial cancellation of the errors and does not change 
the value of $\phi$ by much
relative to the norm $\|\phi\|_\infty$.  A precise
statement follows.

\begin{theorem}
\label{thm:error-compounding-multilinear}
Let $\phi\colon\moon\to\R$ be given such that
$\hat\phi(S)=0$ whenever $|S|\ne d$.
Fix an arbitrary symmetric function
$\delta\colon\moo^d\to\R$ and define
$\Delta\colon\moon\to\R$ by
\begin{align*}
\Delta(x) = \sum_{|S|=d}\hat\phi(S) \delta(x|_S)\,.
\end{align*}
Then
\begin{align*}
\|\Delta\|_\infty \leq 
\frac{d^d}{d!}\, \|\phi\|_\infty \|\hat \delta\|_1\,.
\end{align*}
\end{theorem}

In the above result, $\delta$ should be thought of as
the error in approximating individual characters
$\chi_S,$ whereas $\Delta$ is the cumulative error
so incurred.  The theorem states that the cumulative error
exceeds the norm of $\phi$ and $\delta$ by a factor of
only $\e^d,$ which is substantially smaller than the
factor of $n^{\Omega(d)}$ growth that one could expect
\emph{a priori}.

\begin{proof}[Proof of \expref{Theorem}{thm:error-compounding-multilinear}]
We adopt
the convention that $a^0=1$ for all real $a$. For a
given vector $v\in\zoo^d,$ consider the 
operator $A_v$ that takes 
a function $f\colon\moon\to\R$ into
%
%
%
the
%
function $A_vf\colon\moon\to\R$ where
\begin{align*}
(A_vf)(x) = 
\Exp_{z\in\moo^d} \;\left[ z_1z_2\cdots z_df\left(\frac1d\sum_{i=1}^d
z_ix_1^{v_i}, \dots,
                     \frac1d\sum_{i=1}^d
					 z_ix_n^{v_i}\right)\right].
\end{align*}
This definition is meaningful because we identify $f$
with its multilinear extension to $[-1,1]^n,$ thus making
it possible to evaluate $f$ on inputs with fractional
coordinates (see \expref{Section}{sec:prelim}). 
It is important to note that $A_v$ is a linear
transformation in the vector space $\R^{\moon}$.
This somewhat magical operator is the key to the
proof; the remainder of the proof will provide insight
into how this definition could have been arrived at in the
first place. To start with, 
\begin{align}
\|A_v\phi\|_\infty \leq
\max_{x\in[-1,1]^n}|\phi(x)| = 
\max_{x\in\moon}|\phi(x)| = \|\phi\|_\infty\,,
\label{eqn:A_v-small}
\end{align}
where the second step holds by convexity.  The strategy
of the proof is to express $\Delta$ as a linear
combination of the $A_v\phi$ with small coefficients.
Since the infinity norm of any individual $A_v\phi$
is small, that will give the desired bound on the
infinity norm of $\Delta$. 

To find what suitable coefficients would be, we need to
understand the transformation $A_v$ in terms of the Fourier
spectrum. Since $A_v$ is linear and the nonzero
Fourier coefficients of $\phi$ have order $d,$ it
suffices to determine the action of $A_v$ on the
characters of order $d$. For every $S\subseteq\oneton$
with $|S|=d,$ 
\begin{align}
(A_v\chi_S)(x)
 &= 
\Exp_{z\in\moo^d} \;
\left[z_1z_2\cdots z_d \;\prod_{j\in S}\left(\frac1d\sum_{i=1}^d
z_{i}x_j^{v_i}\right)\right] \nonumber\\
&= 
\Exp_{z\in\moo^d} \;
\left[z_1z_2\cdots z_d \;\Exp_{\tau\colon S\to\{1,\dots,d\}}
\left[\prod_{j\in S}
z_{\tau(j)}x_j^{v_{\tau(j)}}\right]\right] \nonumber\\
&= \Exp_{\tau\colon S\to\{1,\dots,d\}} \;
\left[
\Exp_{z\in\moo^d}\left[z_1z_2\cdots z_d\prod_{j\in S} 
z_{\tau(j)}\right]
\;\prod_{j\in S} x_j^{v_{\tau(j)}}
\right],
\label{eqn:enforce-bijection}
\end{align}
where $\tau$ chosen uniformly at random from all
possible mappings $S\to\{1,2,\dots,d\}$. 
The inner expectation in \eqref{eqn:enforce-bijection} acts like the
indicator random variable for the event that $\tau$ is
a bijection, \ie, it evaluates to $1$ when $\tau$ is a
bijection and vanishes otherwise. As a result,
\begin{align}
(A_v\chi_S)(x)
&= \Prob_\tau [\text{$\tau$ is a bijection}]\;
\Exp_\tau \left[\prod_{j\in S} x_j^{v_{\tau(j)}}
\;\middle|\;  \text{$\tau$ is a bijection}\right]\nonumber\\
 &= 
\frac{d!}{d^d}\;\;
\Exp_{\substack{T\subseteq
S,\\|T|=v_1+\cdots+v_d}}\!\![\chi_T(x)]\,.
\label{eqn:action-on-character}
\end{align}
By the symmetry of $\delta,$ 
\begin{align*}
\delta(x|_S) &= \sum_{k=0}^d \hat\delta(\{1,2,\dots,k\}) \sum_{\substack{T\subseteq
 S,\\|T|=k}}\chi_T(x)\\
 &= \frac{d^d}{d!}\sum_{k=0}^d \hat\delta(\{1,2,\dots,k\}) {d\choose
 k}
 (A_{1^k0^{d-k}}\chi_S)(x)\,,
\end{align*}
where the second step uses
\eqref{eqn:action-on-character}. Taking a weighted sum over $S$ and using the
linearity of $A_v,$ 
\begin{align*}
\sum_{\substack{S\subseteq\oneton\\|S|=d}}
\hat\phi(S)\delta(x|_S) &= 
 \frac{d^d}{d!}\sum_{k=0}^d \hat\delta(\{1,2,\dots,k\}) {d\choose
 k}
 \PARENS{A_{1^k0^{d-k}}\sum_{\substack{S\subseteq\oneton\\|S|=d}}^{\phantom
 {\substack{S\\d}}}
 \hat\phi(S)\chi_S}(x)\,,
\intertext{or equivalently}
\Delta
&= \frac{d^d}{d!} \sum_{k=0}^d \hat\delta(\{1,2,\dots,k\}) {d\choose
 k} A_{1^k0^{d-k}}\phi\,.
\end{align*}
In light of \eqref{eqn:A_v-small}, this representation
gives the sought upper bound on the norm of $\Delta$:
\begin{align*}
\|\Delta\|_\infty \leq
& \frac{d^d}{d!} \sum_{k=0}^d |\hat\delta(\{1,2,\dots,k\})|{d\choose
 k}
 \|\phi\|_\infty
=\frac{d^d}{d!}\, \|\phi\|_\infty \|\hat \delta\|_1\,.
\qedhere
\end{align*}
\end{proof}


\subsection*{Error cancellation with real variables}
We now consider the error cancellation problem in its full
generality. Again, our goal will be to show that
replacing individual characters with suitable
approximants results in moderate cumulative error. 
This time, however, the input variables are no longer restricted
to be Boolean, and can take on arbitrary values in 
\mbox{$[-\sqrt{1+\epsilon},-\sqrt{1-\epsilon}]\cup[\sqrt{1-\epsilon},\sqrt{1+\epsilon}]$}
for $0<\epsilon<1$. This in turn means that the
error term will be given by an infinite series.
Another difference is that the coefficients of the error
series will not converge to zero rapidly enough,
requiring additional ideas to bound the cumulative
error.


\begin{theorem}
\label{thm:error-compounding-Rn}
Let $\phi\colon\moon\to\R$ be given such that
$\hat\phi(S)=0$ whenever $|S|\ne d$.
Let $0<\epsilon<1$ and 
\begin{align*}
X=[-\sqrt{1+\epsilon},-\sqrt{1-\epsilon}]\cup
         [\sqrt{1-\epsilon},\sqrt{1+\epsilon}]\,. 
\end{align*}
Then for every integer $D\geq1,$ there is an
$($explicitly given$)$ polynomial $p\colon\R^d\to\R$
of degree at most $2D+d$ such that
\begin{align*}
P(x) = \sum_{\substack{S\subseteq\oneton\\|S|=d}}\hat\phi(S)p(x|_S)
\end{align*}
obeys 
\begin{align}
   \max_{X^n}\left|\phi(\sign x_1,\dots,\sign
   x_n)-P(x)\right| \leq
(1+\epsilon)^{d/2}\;\frac{d^d}{d!}\;\epsilon^D{D+d\choose D}
d\,\|\phi\|_\infty\,. \hspace{1cm}
   \label{eqn:homog-error-bound}
\end{align}
\end{theorem}


\begin{proof}
As before, we adopt
the notational convention that $a^0=1$ for all real
$a$. We will follow the proof of
\expref{Theorem}{thm:error-compounding-multilinear} as
closely as possible, pointing out key differences as we
go along.

If $\epsilon> D/{(D+d)},$ then the right-hand side of 
\eqref{eqn:homog-error-bound} is at least 
\[
\left(\frac D{D+d}\right)^D{D+d\choose D}
\|\phi\|_\infty\geq \|\phi\|_\infty
\] 
and hence the
theorem holds trivially with $p=0$. In what follows, 
we treat the complementary case 
$\epsilon\leq D/{(D+d)}$. To start with,
\begin{align}
\sum_{\substack{\rule{0mm}{2mm}v\in\mathbb
N^d:\\\rule{0mm}{3mm}v_1+\cdots+v_d\geq D+1}} 
 \epsilon^{v_1+\cdots+v_d}
\prod_{j=1}^d\underbrace{\PARENS{\frac14}^{v_j}{2v_j \choose v_j}}_{\leq 1}
&\leq\sum_{\substack{\rule{0mm}{2mm}v\in\mathbb
N^d:\\\rule{0mm}{3mm}v_1+\cdots+v_d\geq D+1}} 
 \epsilon^{v_1+\cdots+v_d}\nonumber\\
&=\epsilon^{D+1}\sum_{i=0}^\infty{D+d+i\choose D+1+i}\epsilon^i
\nonumber\\
&\leq\epsilon^{D+1}{D+d\choose D+1}\sum_{i=0}^\infty\left(\frac{D+d}{D+1}\right)^i\epsilon^i
\nonumber\\
&\leq\epsilon^{D+1}{D+d\choose D+1}\sum_{i=0}^\infty\left(\frac{D}{D+1}\right)^i
\nonumber\\
&=\epsilon^{D+1}{D+d\choose D}d\,,
\label{eqn:coefficients-converge}
\end{align}
where the fourth line in the derivation uses
$\epsilon\leq D/(D+d)$.
Define $p$ by 
\begin{align*}
p(x_1,\dots,x_d) = 
\sum_{\substack{\rule{0mm}{2mm}v\in\mathbb
N^d:\\\rule{0mm}{3mm}v_1+\cdots+v_d\leq D}} 
\;\; \prod_{j=1}^d\PARENS{-\frac14}^{v_j}{2v_j \choose v_j}
x_j(x_j^2-1)^{v_j}\,.
\end{align*} Analogous to the Boolean setting, we will
define functions to capture the error in approximating
an individual character as well as the 
cumulative error. Let $\delta\colon X^d\to\R$ and
$\Delta\colon X^n\to\R$ be given by 
\begin{align*}
\delta(x_1,\dots,x_d) &=
\sum_{\substack{\rule{0mm}{2mm}v\in\mathbb
N^d:\\\rule{0mm}{3mm}v_1+\cdots+v_d\geq D+1}} 
\;\; \prod_{j=1}^d\PARENS{-\frac14}^{v_j}{2v_j \choose v_j}
x_j(x_j^2-1)^{v_j}\,, \\
\rule{0mm}{9mm}\Delta(x_1,\dots,x_n) &=
\sum_{\substack{S\subseteq\oneton\\|S|=d}}
\hat\phi(S)\delta(x|_S)\,.
\end{align*}
\expref{Lemma}{lem:parity-robust} implies that $\delta$ is
the error incurred in approximating a single
character by $p,$ in other words,
$\delta(x_1,\dots,x_d)= \sign(x_1\cdots
x_d)-p(x_1,\dots,x_d)$. Hence, $\Delta$ captures the cumulative
error:
\begin{align}
\Delta(x) = \phi(\sign x_1,\dots,\sign x_n)-P(x).
\label{eqn:Delta-cumulative}
\end{align}

Recall that our goal is to
place an upper bound on $\|\Delta\|_\infty$. 
For $v\in\mathbb N^d,$ 
consider the operator $A_v$ that
takes a function $f\colon\moon\to\R$ into a function
$A_vf\colon X^n\to\R$ where 
\begin{align*}
(A_vf)(x) = 
\Exp_{z\in\moo^d} \;\left[\rule{0mm}{7mm}
z_1z_2\cdots z_d\;f
\right(
\dots,
\underbrace{\frac1d\sum_{i=1}^d
\frac{z_ix_j(x_j^2-1)^{v_i}}{\epsilon^{v_i}
\sqrt{1+\epsilon}}}_{\text{$j$th
coordinate}}
,\dots\left)\rule{0mm}{7mm}
\rule{0mm}{6mm}\right].
\end{align*}
The $j$th coordinate in this expression is bounded by 
$1$ in absolute value because
$\sqrt{1-\epsilon}<|x_j|<\sqrt{1+\epsilon}$ on $X^n$.
This definition departs from the earlier
one in \expref{Theorem}{thm:error-compounding-multilinear},
where $v$ was restricted to $0/1$ entries. Perhaps the most
essential difference is the presence of scaling
factors in the denominator---it is what ultimately
allows one to bound the cumulative error in the setting
of an infinite series. Note that $A_v$ is a linear
transformation sending $\R^{\moon}$ into $\R^{X^n}$. 
We further have
\begin{align}
\|A_v\phi\|_\infty \leq
\max_{x\in[-1,1]^n}|\phi(x)| = 
\max_{x\in\moon}|\phi(x)| = \|\phi\|_\infty\,,
\label{eqn:A_v-small-Rn}
\end{align}
where the first step uses the fact that $A_v\phi$ has
domain $X^n$ rather than all of $\R^n,$ and the second step holds by convexity. 

We proceed to examine the action of $A_v$ on the
characters of order $d$. Since the definition of $A_v$
is symmetric with respect to the $n$ coordinates, it
suffices to consider
$S=\{1,2,\dots,d\}$:
\begin{align*}
(A_v\chi_{\{1,\dots,d\}})(x)
 &=
\Exp_{z\in\moo^d} \;
\left[z_1z_2\cdots z_d \prod_{j=1}^d\left(
\frac1d\sum_{i=1}^d
\frac{z_ix_j(x_j^2-1)^{v_i}}
{\epsilon^{v_i}
\sqrt{1+\epsilon}
}
\right)
\right]
\\
&=
\frac1{(1+\epsilon)^{d/2}}\;
 \Exp_{\tau} 
\left[\Exp_{z}\left[\prod_{j=1}^d z_j
z_{\tau(j)}\right]\;\prod_{j=1}^d
\frac{x_j(x_j^2-1)^{v_{\tau(j)}}}
     {\epsilon^{v_{\tau(j)}}} 
\right],
\end{align*}
where the first expectation is taken over a uniformly
random mapping $\tau\colon \{1,2,\dots,d\}\to\{1,2,\dots,d\}$. 
The expectation over $z$ acts like the
indicator random variable for the event that $\tau$ is a bijection,
\ie, it evaluates to $1$ when $\tau$ is a bijection and
vanishes otherwise.  Thus,
\begin{align}
(A_v\chi_{\{1,\dots,d\}})(x)
&=
\frac1{(1+\epsilon)^{d/2}}\;
\Prob_\tau[\text{$\tau$ is a bijection}] \;
\Exp_{\tau} 
\left[\prod_{j=1}^d
\frac{x_j(x_j^2-1)^{v_{\tau(j)}}}
     {\epsilon^{v_{\tau(j)}}} \;
\middle| \; \text{$\tau$ is a bijection}\right]\nonumber\\
&=
\frac1{(1+\epsilon)^{d/2}\epsilon^{v_1+\cdots+v_d}}
\cdot \frac{d!}{d^d}\cdot
\Exp_{\sigma\in S_d} 
\left[\,\prod_{j=1}^d x_j(x_j^2-1)^{v_{\sigma(j)}}\right]\,.
\label{eqn:action-on-character-Rn}
\end{align}
Now, consider the operator
\begin{align*}
A= 
(1+\epsilon)^{d/2}\;\frac{d^d}{d!}\sum_{\substack{\rule{0mm}{2mm}v\in\mathbb
N^d:\\\rule{0mm}{3mm}v_1+\cdots+v_d\geq D+1}} 
 \BRACES{\epsilon^{v_1+\cdots+v_d}
\prod_{j=1}^d\PARENS{-\frac14}^{v_j}{2v_j \choose v_j}}
  \; A_v\,.
\end{align*}
This operator is well-defined because by 
\eqref{eqn:coefficients-converge},
the series in question converges absolutely.
By \eqref{eqn:action-on-character-Rn}, the symmetry
of $\delta,$ and linearity,
$
(A\chi_{\{1,\dots,d\}})(x)=\delta(x_1,\dots,x_d)\,.
$
Since the definition of $A$ is symmetric with respect to
the $n$ coordinates, we conclude that
\begin{align*}
(A\chi_S)(x) = \delta(x|_S)
\end{align*}
for all subsets $S\subseteq\oneton$ of cardinality $d$.
As an immediate consequence, 
\begin{align*}
\Delta = \sum_{\substack{S\subseteq\oneton\\|S|=d}}
	\hat\phi(S)\cdot (A\chi_S)
       = A
	   \PARENS{
	   \sum_{\substack{S\subseteq\oneton\\|S|=d}}
	       ^{\phantom{\substack{|\\S\subseteq}}}
	   \hat\phi(S)\chi_S}
       = A\phi\,,
\end{align*}
where the second step uses the linearity of $A$. In 
particular,
\begin{align*}
\|\Delta\|_\infty &=\|A\phi\|_\infty \\
&\leq (1+\epsilon)^{d/2}\;\frac{d^d}{d!}\sum_{\substack{\rule{0mm}{2mm}v\in\mathbb
N^d:\\\rule{0mm}{3mm}v_1+\cdots+v_d\geq D+1}} 
 \epsilon^{v_1+\cdots+v_d}
  \; \|A_v\phi\|_\infty\prod_{j=1}^d
  \PARENS{\frac14}^{v_j}{2v_j\choose v_j}\\
&\rule{0mm}{9mm}\leq
(1+\epsilon)^{d/2}\;\frac{d^d}{d!}\;\epsilon^{D+1}{D+d\choose D}
d\,\|\phi\|_\infty\,,
\end{align*}
where the final step follows by
\eqref{eqn:coefficients-converge} and
\eqref{eqn:A_v-small-Rn}.
In light of \eqref{eqn:Delta-cumulative}, the proof is
complete.
\end{proof}




\section{Main result}
\label{sec:main}


We are now in a position to prove the main result of this
paper, which states that every bounded real polynomial 
can be made robust with only a constant-factor increase
in degree. Recall that we have already proved this fact for
homogeneous polynomials (see
Theorems~\ref{thm:error-compounding-multilinear}
and~\ref{thm:error-compounding-Rn}). It remains to
remove the homogeneity assumption, which we will do using
the technique of \expref{Section}{sec:homogeneous}.
For the purposes of exposition, we will first show how
to remove the homogeneity assumption in the much
simpler context of
\expref{Theorem}{thm:error-compounding-multilinear}.
Essentially the same technique will then allow us to
prove the main result.

\begin{theorem}
\label{thm:cumulative-error-nonhomog-multilinear}
Let $\phi\colon\moon\to\R$ be given, $\deg\phi=d$.
Fix symmetric functions $\delta_i\colon\moo^i\to\R,$
$i=0,1,2,\dots,d,$
and define $\Delta\colon\moon\to\R$ by
\begin{align*}
\Delta(x)=\sum_{|S|\leq d}
\hat\phi(S)\delta_{|S|}(x|_S)\,.
\end{align*}
Then
\begin{align*}
\|\Delta\|_\infty \leq (4\e)^d\, \|\phi\|_\infty
\,\sum_{i=0}^d \|\hat\delta_i\|_1\,.
\end{align*}
\end{theorem}

The functions $\delta_0,\delta_1,\dots,\delta_d$ in
this result are to be thought of as the errors in the approximation of 
characters of orders $0,1,\dots,d,$ respectively, and
$\Delta$ is the cumulative error so incurred.
As the theorem shows, the
cumulative error exceeds the norms of the functions
involved by a factor of only $2^{O(d)},$ which is
independent of the number of variables.

\begin{proof}
[Proof of \expref{Theorem}{thm:cumulative-error-nonhomog-multilinear}]
We have 
\begin{align}
\|\Delta\|_\infty \leq
\sum_{i=0}^d\|\Delta_i\|_\infty\,,
\label{eqn:Delta-Delta_i}
\end{align}
where $\Delta_i\colon\moon\to\R$ is given by
$\Delta_i(x) = \sum_{|S|=i}\hat\phi(S)\delta_i(x|_S)$.
For $i=0,1,\dots,d,$ consider 
$\phi_i=\sum_{|S|=i}\hat\phi(S)\chi_S,$ the
degree-$i$ homogeneous part of $\phi$.
By \expref{Theorem}{thm:error-compounding-multilinear},
\begin{align}
\|\Delta_i\|_\infty \leq \e^i \|\phi_i\|_\infty\:
\|\hat\delta_i\|_1\,.
\end{align}
By \expref{Theorem}{thm:rebuild-from-levels},
\begin{align}
\|\phi_i\|_\infty \leq 4^d \|\phi\|_\infty\,, \qquad
i=0,1,\dots,d\,.
\label{eqn:phi_i-phi}
\end{align}
Combining
\eqref{eqn:Delta-Delta_i}--\eqref{eqn:phi_i-phi}
completes the proof.
\end{proof}


We will now apply a similar argument in the setting of
real variables. 

%
\begin{theorem}[Main Theorem]
Let $0<\epsilon<1$ and $\phi\colon\moon\to[-1,1]$ be given,
$\deg\phi=d$. 
Then for each $\delta>0,$ there is a
polynomial $P$ of degree
\[
O\left(\frac1{1-\epsilon}
d+\frac1{1-\epsilon} \log \frac1\delta\right)
\]
such that 
\begin{align}
   \max_{x\in X^n}\left|\phi(\sign x_1,\dots,\sign
   x_n)-P(x)\right| <\delta\,,
   \label{eqn:P-main-technical}
\end{align}
where $X=[-1-\epsilon,-1+\epsilon]\cup
         [1-\epsilon,1+\epsilon]$.
Furthermore, $P$ is given explicitly in terms of the
Fourier spectrum of $\phi$. 
\end{theorem}
%

Letting $\epsilon=1/3$ in this result settles \expref{Theorem}{thm:main} in the introduction. 

\begin{proof}
We first consider $0\leq\epsilon\leq1/10,$ in which
case 
\begin{align}
X\subset\left[-\sqrt{1+\frac14},-\sqrt{1-\frac14}\right]\cup
\left[\sqrt{1-\frac14},\sqrt{1+\frac14}\right]\,.
\label{eqn:X}
\end{align}
Let $D=D(d,\delta)\geq1$ be a parameter to be chosen later.
For $i=0,1,\dots,d,$ consider 
$\phi_i=\sum_{|S|=i}\hat\phi(S)\chi_S,$ the
degree-$i$ homogeneous part of $\phi$.
By \expref{Theorem}{thm:rebuild-from-levels},
\begin{align}
\label{eqn:phi_i_bounded}
\|\phi_i\|_\infty \leq 4^d, \qquad
i=0,1,\dots,d\,.
\end{align}
In light of \eqref{eqn:X}, \expref{Theorem}{thm:error-compounding-Rn} gives explicit
polynomials $p_i\colon\R^i\to\R,$
$i=0,1,2,\dots,d,$ each of degree at most $2D+d,$ such
that
\begin{align*}
\max_{X^n}\Biggl|\phi_i(\sign x_1,\dots,\sign x_n)
 - \sum_{\substack{S\subseteq\oneton\\|S|=i}}\hat\phi(S)p_i(x|_S)
 \Biggr| \leq
\frac{K^d}{2^D} \, \|\phi_i\|_\infty
\end{align*}
for some absolute constant $K>1$. 
Letting 
\begin{align*}
P(x) = \sum_{|S|\leq d}\hat\phi(S)p_{|S|}(x|_S)\,, 
\end{align*}
we infer that 
\begin{align*}
\max_{X^n}\left|\phi(\sign x_1,\dots,\sign x_n)
 - P(x)\right|
\leq
\frac{K^d}{2^D} \sum_{i=0}^d\|\phi_i\|_\infty
\leq 
\frac{(d+1)(4K)^d}{2^D}\,,
\end{align*}
where the last step uses \eqref{eqn:phi_i_bounded}.
Therefore, \eqref{eqn:P-main-technical} holds with 
$D=O(d+\log 1/\delta)$. 

To handle the case $1/10<\epsilon<1,$ basic approximation
theory~\cite{rivlin-book} gives an explicit univariate
polynomial $r$ of degree $O(1/{(1-\epsilon)})$ that sends
$[-1-\epsilon,-1+\epsilon]\to[-11/10,-9/10]$ and 
$[1-\epsilon,1+\epsilon]\to[9/10,11/10]$. In particular, we have
$\left|\phi(\sign x_1,\dots,\sign
x_n)-P(r(x_1),\dots,r(x_n))\right|<\delta$ everywhere
on $X^n,$ where $P$ is the approximant constructed in the
previous paragraph.
\end{proof}

%
%
%
%
%
\providecommand{\bibhead}[1]{}
%
\expandafter\ifx\csname pdfbookmark\endcsname\relax%
  \providecommand{\tocrefpdfbookmark}{}
\else\providecommand{\tocrefpdfbookmark}{%
   \hypertarget{tocreferences}{}%
   \pdfbookmark[1]{References}{tocreferences}}%
\fi

\tocrefpdfbookmark
\begin{thebibliography}{10}

\bibitem{aaronson04sdpt-for-search}\bibhead{aaronson04sdpt-for-search}
{\sc Scott Aaronson}: Limitations of quantum advice and one-way communication.
\newblock {\em Theory of Computing}, 1(1):1--28, 2005.
\newblock Preliminary version in
  \href{http://dx.doi.org/10.1109/CCC.2004.1313854}{CCC'04}.
\newblock [\epfmtdoi{10.4086/toc.2005.v001a001}]

\bibitem{aaronson-shi04distinctness}\bibhead{aaronson-shi04distinctness}
{\sc Scott Aaronson and Yaoyun Shi}: Quantum lower bounds for the collision and
  the element distinctness problems.
\newblock {\em J. ACM}, 51(4):595--605, 2004.
\newblock Preliminary version in
  \href{http://dx.doi.org/10.1145/509907.509999}{STOC'02}.
\newblock [\epfmtdoi{10.1145/1008731.1008735}]

\bibitem{aspnes91voting}\bibhead{aspnes91voting}
{\sc James Aspnes, Richard Beigel, Merrick~L. Furst, and Steven Rudich}: The
  expressive power of voting polynomials.
\newblock {\em Combinatorica}, 14(2):135--148, 1994.
\newblock Preliminary version in
  \href{http://dx.doi.org/10.1145/103418.103461}{STOC'91}.
\newblock [\epfmtdoi{10.1007/BF01215346}]

\bibitem{beals-et-al01quantum-by-polynomials}\bibhead{beals-et-al01quantum-by-polynomials}
{\sc Robert Beals, Harry Buhrman, Richard Cleve, Michele Mosca, and Ronald~de
  Wolf}: Quantum lower bounds by polynomials.
\newblock {\em J. ACM}, 48(4):778--797, 2001.
\newblock Preliminary version in
  \href{http://dx.doi.org/10.1109/SFCS.1998.743485}{FOCS'98}.
\newblock [\epfmtdoi{10.1145/502090.502097}]

\bibitem{beame-huyn-ngoc09multiparty-focs}\bibhead{beame-huyn-ngoc09multiparty-focs}
{\sc Paul Beame and Dang-Trinh {Huynh-Ngoc}}: Multiparty communication
  complexity and threshold circuit size of {$\mathsf{AC}^0$}.
\newblock {\em SIAM Journal on Computing}, 41(3):484--518, 2012.
\newblock Preliminary version in
  \href{http://dx.doi.org/10.1109/FOCS.2009.12}{FOCS'09}.
\newblock [\epfmtdoi{10.1137/100792779}]

\bibitem{beigel91rational}\bibhead{beigel91rational}
{\sc Richard Beigel, Nick Reingold, and Daniel~A. Spielman}: {$\mathsf{PP}$} is
  closed under intersection.
\newblock {\em J. Comput. System Sci.}, 50(2):191--202, 1995.
\newblock Preliminary version in
  \href{http://dx.doi.org/10.1145/103418.103426}{STOC'91}.
\newblock [\epfmtdoi{10.1006/jcss.1995.1017}]

\bibitem{braverman-rao11noisy-communication}\bibhead{braverman-rao11noisy-communication}
{\sc Mark Braverman and Anup Rao}: Towards coding for maximum errors in
  interactive communication.
\newblock In {\em Proc. 43rd STOC}, pp. 159--166. ACM Press, 2011.
\newblock [\epfmtdoi{10.1145/1993636.1993659}]

\bibitem{bcw98quantum}\bibhead{bcw98quantum}
{\sc Harry Buhrman, Richard Cleve, and Avi Wigderson}: Quantum vs. classical
  communication and computation.
\newblock In {\em Proc. 30th STOC}, pp. 63--68. ACM Press, 1998.
\newblock [\epfmtdoi{10.1145/276698.276713}]

\bibitem{buhrman-et-al99small-error}\bibhead{buhrman-et-al99small-error}
{\sc Harry Buhrman, Richard Cleve, Ronald~de Wolf, and Christof Zalka}: Bounds
  for small-error and zero-error quantum algorithms.
\newblock In {\em Proc. 40th FOCS}, pp. 358--368. IEEE Comp. Soc. Press, 1999.
\newblock [\epfmtdoi{10.1109/SFFCS.1999.814607}]

\bibitem{robust}\bibhead{robust}
{\sc Harry Buhrman, Ilan Newman, Hein R{\"o}hrig, and Ronald~de Wolf}: Robust
  polynomials and quantum algorithms.
\newblock {\em Theory Comput. Syst.}, 40(4):379--395, 2007.
\newblock Preliminary version in
  \href{http://dx.doi.org/10.1007/978-3-540-31856-9_49}{STACS'05}.
\newblock [\epfmtdoi{10.1007/s00224-006-1313-z}]

\bibitem{buhrman07pp-upp}\bibhead{buhrman07pp-upp}
{\sc Harry Buhrman, Nikolai~K. Vereshchagin, and Ronald de~Wolf}: On
  computation and communication with small bias.
\newblock In {\em Proc. 22nd IEEE Conf. on Computational Complexity (CCC'07)},
  pp. 24--32. IEEE Comp. Soc. Press, 2007.
\newblock (See also
  \href{http://dx.doi.org/10.1007/978-3-540-74871-7_1}{SAGA'07}.).
\newblock [\epfmtdoi{10.1109/CCC.2007.18}]

\bibitem{buhrman-dewolf01polynomials}\bibhead{buhrman-dewolf01polynomials}
{\sc Harry Buhrman and Ronald~de Wolf}: Communication complexity lower bounds
  by polynomials.
\newblock In {\em Proc. 16th IEEE Conf. on Computational Complexity (CCC'01)},
  pp. 120--130. IEEE Comp. Soc. Press, 2001.
\newblock [\epfmtdoi{10.1109/CCC.2001.933879}]

\bibitem{cheney-book}\bibhead{cheney-book}
{\sc Elliott~W. Cheney}: {\em Introduction to Approximation Theory}.
\newblock Chelsea Publishing, New York, 2nd edition, 1982.

\bibitem{DKMR08noisy-parity-broadcast-networks}\bibhead{DKMR08noisy-parity-broadcast-networks}
{\sc Chinmoy Dutta, Yashodhan Kanoria, D.~Manjunath, and Jaikumar
  Radhakrishnan}: A tight lower bound for parity in noisy communication
  networks.
\newblock In {\em Proc. 19th Ann. ACM-SIAM Symp. on Discrete Algorithms
  (SODA'08)}, pp. 1056--1065. ACM Press, 2008.
\newblock [\epfmt{acm}{1347198}]

\bibitem{dutta-radhakrishnan08noisy-broadcast}\bibhead{dutta-radhakrishnan08noisy-broadcast}
{\sc Chinmoy Dutta and Jaikumar Radhakrishnan}: Lower bounds for noisy wireless
  networks using sampling algorithms.
\newblock In {\em Proc. 49th FOCS}, pp. 394--402. IEEE Comp. Soc. Press, 2008.
\newblock [\epfmtdoi{10.1109/FOCS.2008.72}]

\bibitem{eremenko07sign}\bibhead{eremenko07sign}
{\sc Alexandre Eremenko and Peter Yuditskii}: Uniform approximation of
  $\mathrm{sgn}~x$ by polynomials and entire functions.
\newblock {\em J. d'Analyse Math{\'e}matique}, 101(1):313--324, 2007.
\newblock [\epfmtdoi{10.1007/s11854-007-0011-3}]

\bibitem{evans-pippenger96noisy-decision-trees}\bibhead{evans-pippenger96noisy-decision-trees}
{\sc William~S. Evans and Nicholas Pippenger}: Average-case lower bounds for
  noisy {B}oolean decision trees.
\newblock {\em SIAM J. Comput.}, 28(2):433--446, 1998.
\newblock Preliminary version in
  \href{http://dx.doi.org/10.1145/237814.238013}{STOC'96}.
\newblock [\epfmtdoi{10.1137/S0097539796310102}]

\bibitem{evans-schulman99noisy-circuits}\bibhead{evans-schulman99noisy-circuits}
{\sc William~S. Evans and Leonard~J. Schulman}: Signal propagation and noisy
  circuits.
\newblock {\em IEEE Trans. Inform. Theory}, 45(7):2367--2373, 1999.
\newblock Preliminary version in
  \href{http://dx.doi.org/10.1109/SFCS.1993.366827}{FOCS'93}.
\newblock [\epfmtdoi{10.1109/18.796377}]

\bibitem{feige92finite-random-functions}\bibhead{feige92finite-random-functions}
{\sc Uriel Feige}: On the complexity of finite random functions.
\newblock {\em Inform. Process. Lett.}, 44(6):295--296, 1992.
\newblock [\epfmtdoi{10.1016/0020-0190(92)90102-2}]

\bibitem{feige-kilian00fiding-OR-noisy-broadcast}\bibhead{feige-kilian00fiding-OR-noisy-broadcast}
{\sc Uriel Feige and Joe Kilian}: Finding {OR} in a noisy broadcast network.
\newblock {\em Inform. Process. Lett.}, 73(1-2):69--75, 2000.
\newblock [\epfmtdoi{10.1016/S0020-0190(00)00002-8}]

\bibitem{FRPU94noise}\bibhead{FRPU94noise}
{\sc Uriel Feige, Prabhakar Raghavan, David Peleg, and Eli Upfal}: Computing
  with noisy information.
\newblock {\em SIAM J. Comput.}, 23(5):1001--1018, 1994.
\newblock Preliminary version in
  \href{http://dx.doi.org/10.1145/100216.100230}{STOC'90}.
\newblock [\epfmtdoi{10.1137/S0097539791195877}]

\bibitem{gacs-gal94noisy-circuits}\bibhead{gacs-gal94noisy-circuits}
{\sc P{\'e}ter G{\'a}cs and Anna G{\'a}l}: Lower bounds for the complexity of
  reliable {B}oolean circuits with noisy gates.
\newblock {\em IEEE Trans. Inform. Theory}, 40(2):579--583, 1994.
\newblock Preliminary version in
  \href{http://dx.doi.org/10.1109/SFCS.1991.185424}{FOCS'91}.
\newblock [\epfmtdoi{10.1109/18.312190}]

\bibitem{gallager88noisy-broadcast}\bibhead{gallager88noisy-broadcast}
{\sc Robert~G. Gallager}: Finding parity in a simple broadcast network.
\newblock {\em IEEE Trans. Inform. Theory}, 34(2):176--180, 1988.
\newblock [\epfmtdoi{10.1109/18.2626}]

\bibitem{GMS11noisy-communication}\bibhead{GMS11noisy-communication}
{\sc Ran Gelles, Ankur Moitra, and Amit Sahai}: Efficient and explicit coding
  for interactive communication.
\newblock In {\em Proc. 52nd FOCS}, pp. 768--777. IEEE Comp. Soc. Press, 2011.
\newblock [\epfmtdoi{10.1109/FOCS.2011.51}]

\bibitem{GKS05noisy-broadcast}\bibhead{GKS05noisy-broadcast}
{\sc Navin Goyal, Guy Kindler, and Michael~E. Saks}: Lower bounds for the noisy
  broadcast problem.
\newblock {\em SIAM J. Comput.}, 37(6):1806--1841, 2008.
\newblock Preliminary version in
  \href{http://dx.doi.org/10.1109/SFCS.2005.48}{FOCS'05}.
\newblock [\epfmtdoi{10.1137/060654864}]

\bibitem{concrete-math}\bibhead{concrete-math}
{\sc Ronald~L. Graham, Donald~E. Knuth, and Oren Patashnik}: {\em Concrete
  Mathematics: A Foundation for Computer Science}.
\newblock Addison-Wesley, Boston, Mass., USA, 2nd edition, 1994.

\bibitem{hoyer-mosca-dewolf03and-or-tree}\bibhead{hoyer-mosca-dewolf03and-or-tree}
{\sc Peter H{\o}yer, Michele Mosca, and Ronald~de Wolf}: Quantum search on
  bounded-error inputs.
\newblock In {\em Proc. 30th Internat. Colloq. on Automata, Languages and
  Programming (ICALP'03)}, pp. 291--299. Springer, 2003.
\newblock [\epfmtdoi{10.1007/3-540-45061-0\_25}]

\bibitem{KKMS}\bibhead{KKMS}
{\sc Adam~Tauman Kalai, Adam~R. Klivans, Yishay Mansour, and Rocco~A.
  Servedio}: Agnostically learning halfspaces.
\newblock {\em SIAM J. Comput.}, 37(6):1777--1805, 2008.
\newblock Preliminary version in
  \href{http://dx.doi.org/10.1109/SFCS.2005.13}{FOCS'05}.
\newblock [\epfmtdoi{10.1137/060649057}]

\bibitem{klauck07product-thms}\bibhead{klauck07product-thms}
{\sc Hartmut Klauck, Robert {\v S}palek, and Ronald~de Wolf}: Quantum and
  classical strong direct product theorems and optimal time-space tradeoffs.
\newblock {\em SIAM J. Comput.}, 36(5):1472--1493, 2007.
\newblock Preliminary version in
  \href{http://dx.doi.org/10.1109/FOCS.2004.52}{FOCS'04}.
\newblock [\epfmtdoi{10.1137/05063235X}]

\bibitem{KLM94noisy-circuits}\bibhead{KLM94noisy-circuits}
{\sc Daniel~J. Kleitman, Frank~Thomson Leighton, and Yuan Ma}: On the design of
  reliable {B}oolean circuits that contain partially unreliable gates.
\newblock {\em J. Comput. System Sci.}, 55(3):385--401, 1997.
\newblock Preliminary version in
  \href{http://dx.doi.org/10.1109/SFCS.1994.365682}{FOCS'94}.
\newblock [\epfmtdoi{10.1006/jcss.1997.1531}]

\bibitem{KS01dnf}\bibhead{KS01dnf}
{\sc Adam~R. Klivans and Rocco~A. Servedio}: Learning {DNF} in time {$2^{\tilde
  O(n^{1/3})}$}.
\newblock {\em J. Comput. System Sci.}, 68(2):303--318, 2004.
\newblock Preliminary version in
  \href{http://dx.doi.org/10.1145/380752.380809}{STOC'01}.
\newblock [\epfmtdoi{10.1016/j.jcss.2003.07.007}]

\bibitem{mlj07sq}\bibhead{mlj07sq}
{\sc Adam~R. Klivans and Alexander~A. Sherstov}: Unconditional lower bounds for
  learning intersections of halfspaces.
\newblock {\em Machine Learning}, 69(2-3):97--114, 2007.
\newblock Preliminary version in
  \href{http://dx.doi.org/10.1007/11776420_26}{COLT'06}.
\newblock [\epfmtdoi{10.1007/s10994-007-5010-1}]

\bibitem{colt07rankeps}\bibhead{colt07rankeps}
{\sc Adam~R. Klivans and Alexander~A. Sherstov}: Lower bounds for agnostic
  learning via approximate rank.
\newblock {\em Comput. Complexity}, 19(4):581--604, 2010.
\newblock Preliminary version in
  \href{http://dx.doi.org/10.1007/978-3-540-72927-3_30}{COLT'07}.
\newblock [\epfmtdoi{10.1007/s00037-010-0296-y}]

\bibitem{krause94depth2mod}\bibhead{krause94depth2mod}
{\sc Matthias Krause and Pavel Pudl{\'a}k}: On the computational power of
  depth-2 circuits with threshold and modulo gates.
\newblock {\em Theoret. Comput. Sci.}, 174(1-2):137--156, 1997.
\newblock Preliminary version in
  \href{http://dx.doi.org/10.1145/195058.195103}{STOC'94}.
\newblock [\epfmtdoi{10.1016/S0304-3975(96)00019-9}]

\bibitem{KP98threshold}\bibhead{KP98threshold}
{\sc Matthias Krause and Pavel Pudl{\'a}k}: Computing {B}oolean functions by
  polynomials and threshold circuits.
\newblock {\em Comput. Complexity}, 7(4):346--370, 1998.
\newblock Preliminary version in
  \href{http://dx.doi.org/10.1109/SFCS.1995.492670}{FOCS'95}.
\newblock [\epfmtdoi{10.1007/s000370050015}]

\bibitem{kushilevitz-mansour98noisy-radio}\bibhead{kushilevitz-mansour98noisy-radio}
{\sc Eyal Kushilevitz and Yishay Mansour}: Computation in noisy radio networks.
\newblock {\em SIAM J. Discrete Math.}, 19(1):96--108, 2005.
\newblock Preliminary version in
  \href{http://dl.acm.org/citation.cfm?id=314709}{SODA'98}.
\newblock [\epfmtdoi{10.1137/S0895480103434063}]

\bibitem{markov-v-a1982least-dev}\bibhead{markov-v-a1982least-dev}
{\sc Vladimir~A. Markov}: On functions of least deviation from zero in a given
  interval.
\newblock Russian Academy of Sciences, St. Petersburg, 1892.
\newblock In Russian.

\bibitem{newman04noisy-broadcast}\bibhead{newman04noisy-broadcast}
{\sc Ilan Newman}: Computing in fault tolerant broadcast networks and noisy
  decision trees.
\newblock {\em Random Structures {\&} Algorithms}, 34(4):478--501, 2009.
\newblock Preliminary version in
  \href{http://dx.doi.org/10.1109/CCC.2004.1313813}{CCC'04}.
\newblock [\epfmtdoi{10.1002/rsa.20240}]

\bibitem{paturi-saks94rational}\bibhead{paturi-saks94rational}
{\sc Ramamohan Paturi and Michael~E. Saks}: Approximating threshold circuits by
  rational functions.
\newblock {\em Inform. and Comput.}, 112(2):257--272, 1994.
\newblock Preliminary version in
  \href{http://dx.doi.org/10.1109/FSCS.1990.89559}{FOCS'90}.
\newblock [\epfmtdoi{10.1006/inco.1994.1059}]

\bibitem{pippenger85noisy}\bibhead{pippenger85noisy}
{\sc Nicholas Pippenger}: On networks of noisy gates.
\newblock In {\em Proc. 26th FOCS}, pp. 30--38. IEEE Comp. Soc. Press, 1985.
\newblock [\epfmtdoi{10.1109/SFCS.1985.41}]

\bibitem{razborov02quantum}\bibhead{razborov02quantum}
{\sc Alexander~A. Razborov}: Quantum communication complexity of symmetric
  predicates.
\newblock {\em Izvestiya: Mathematics}, 67(1):145--159, 2003.
\newblock [\epfmtdoi{10.1070/IM2003v067n01ABEH000422}]

\bibitem{RS07dc-dnf}\bibhead{RS07dc-dnf}
{\sc Alexander~A. Razborov and Alexander~A. Sherstov}: The sign-rank of
  {$\mathsf{AC}^0$}.
\newblock {\em SIAM J. Comput.}, 39(5):1833--1855, 2010.
\newblock Preliminary version in
  \href{http://dx.doi.org/10.1109/FOCS.2008.19}{FOCS'08}.
\newblock [\epfmtdoi{10.1137/080744037}]

\bibitem{reischuk-schmeltz91noisy-circuits-decision-trees}\bibhead{reischuk-schmeltz91noisy-circuits-decision-trees}
{\sc R{\"u}diger Reischuk and Bernd Schmeltz}: Reliable computation with noisy
  circuits and decision trees--a general $n\log n$ lower bound.
\newblock In {\em Proc. 32nd FOCS}, pp. 602--611. IEEE Comp. Soc. Press, 1991.
\newblock [\epfmtdoi{10.1109/SFCS.1991.185425}]

\bibitem{rivlin-book}\bibhead{rivlin-book}
{\sc Theodore~J. Rivlin}: {\em An Introduction to the Approximation of
  Functions}.
\newblock Dover Publications, New York, 1981.

\bibitem{schulman92noisy-communication2}\bibhead{schulman92noisy-communication2}
{\sc Leonard~J. Schulman}: Communication on noisy channels: A coding theorem
  for computation.
\newblock In {\em Proc. 33rd FOCS}, pp. 724--733. IEEE Comp. Soc. Press, 1992.
\newblock [\epfmtdoi{10.1109/SFCS.1992.267778}]

\bibitem{schulman96interactive-communication}\bibhead{schulman96interactive-communication}
{\sc Leonard~J. Schulman}: Coding for interactive communication.
\newblock {\em IEEE Trans. Inform. Theory}, 42(6):1745--1756, 1996.
\newblock Preliminary versions in
  \href{http://dx.doi.org/10.1109/SFCS.1992.267778}{FOCS'92} and
  \href{http://dx.doi.org/10.1145/167088.167279}{STOC'93}.
\newblock [\epfmtdoi{10.1109/18.556671}]

\bibitem{dual-survey}\bibhead{dual-survey}
{\sc Alexander~A. Sherstov}: Communication lower bounds using dual polynomials.
\newblock {\em Bulletin of the {EATCS}}, 95:59--93, 2008.
\newblock \href{http://www.eatcs.org/images/bulletin/beatcs95.pdf}{EATCS}.

\bibitem{sherstov09hshs}\bibhead{sherstov09hshs}
{\sc Alexander~A. Sherstov}: The intersection of two halfspaces has high
  threshold degree.
\newblock In {\em Proc. 50th FOCS}, pp. 343--362. IEEE Comp. Soc. Press, 2009.
\newblock [\epfmtdoi{10.1109/FOCS.2009.18}]

\bibitem{sherstov07ac-majmaj}\bibhead{sherstov07ac-majmaj}
{\sc Alexander~A. Sherstov}: Separating {$\mathsf{AC}^0$} from depth-2 majority
  circuits.
\newblock {\em SIAM J. Comput.}, 38(6):2113--2129, 2009.
\newblock Preliminary version in
  \href{http://dx.doi.org/10.1145/1250790.1250834}{STOC'07}.
\newblock [\epfmtdoi{10.1137/08071421X}]

\bibitem{sherstov09opthshs}\bibhead{sherstov09opthshs}
{\sc Alexander~A. Sherstov}: Optimal bounds for sign-representing the
  intersection of two halfspaces by polynomials.
\newblock In {\em Proc. 42nd STOC}, pp. 523--532. ACM Press, 2010.
\newblock [\epfmtdoi{10.1145/1806689.1806762}]

\bibitem{siu-roy-kailath94rational}\bibhead{siu-roy-kailath94rational}
{\sc Kai-Yeung Siu, Vwani~P. Roychowdhury, and Thomas Kailath}: Rational
  approximation techniques for analysis of neural networks.
\newblock {\em IEEE Trans. Inform. Theory}, 40(2):455--466, 1994.
\newblock Preliminary version in
  \href{http://books.nips.cc/nips05.html}{NIPS'92}.
\newblock [\epfmtdoi{10.1109/18.312168}]

\bibitem{spielman96fault-tolerant-parallel}\bibhead{spielman96fault-tolerant-parallel}
{\sc Daniel~A. Spielman}: Highly fault-tolerant parallel computation.
\newblock In {\em Proc. 37th FOCS}, pp. 154--163. IEEE Comp. Soc. Press, 1996.
\newblock [\epfmtdoi{10.1109/SFCS.1996.548474}]

\bibitem{szegedy-chen02noisy-boolean-fns}\bibhead{szegedy-chen02noisy-boolean-fns}
{\sc Mario Szegedy and Xiaomin Chen}: Computing {B}oolean functions from
  multiple faulty copies of input bits.
\newblock {\em Theoret. Comput. Sci.}, 321(1):149--170, 2004.
\newblock Preliminary version in
  \href{http://dx.doi.org/10.1007/3-540-45995-2_47}{LATIN'02}.
\newblock [\epfmtdoi{10.1016/j.tcs.2003.07.001}]

\bibitem{tt99DNF-incl-excl}\bibhead{tt99DNF-incl-excl}
{\sc Jun Tarui and Tatsuie Tsukiji}: Learning {DNF} by approximating
  inclusion-exclusion formulae.
\newblock In {\em Proc. 14th IEEE Conf. on Computational Complexity (CCC'99)},
  pp. 215--221. IEEE Comp. Soc. Press, 1999.
\newblock [\epfmtdoi{10.1109/CCC.1999.766279}]

\end{thebibliography}

